\section{Artifact Availability and Limits}

The artifact repository contains the \tool Python package, adapter scripts,
metric runners, generated CSV/Markdown/\LaTeX{} tables, provenance logs,
BibTeX/source-audit notes, an MIT license, and a reviewer quickstart. The
reviewer entry point is the tag
\texttt{audit-sid-cikm-resource-v0.1}, with a clean-checkout verifier for the
published tables and figure; the anonymous review artifact is available at
\url{https://anonymous.4open.science/r/sidinspector-9BB2}. Large local artifacts are kept out of git and
described through manifests so reviewers can separate public reproducibility
assets from local caches.

The resource is intended to be used in three modes. A downstream recommender
researcher can use the diagnostics as a pre-training triage step: artifacts
with missing items, inconsistent depth, severe aliasing, or weak
collaborative-prefix recovery can be flagged for inspection before generator
training consumes GPU time. A method author can start from the minimal adapter
template, emit the normalized SID-assignment table, then receive D1--D5 reports
without changing their training loop. The artifact also includes runtime notes:
D1/D2/D4/D5 are mapping-level preflight checks, while D3 is the dominant
interaction-neighborhood step. A paper reviewer can run the clean-check verifier
and compare published CSVs with the PDF tables. These signals do not replace
Recall@K or NDCG; they make tokenizer artifacts inspectable before a full
benchmark or generator-training run.

The evidence supports a resource-interface claim, not broad method coverage.
The GRID Musical row is a controlled feature-text export rather than an
end-to-end TIGER or GRID reproduction; ReSID is the bounded Musical export;
RQ-min is a local reference adapter; the Sports balanced GAOQ path was stopped
after constrained k-means became CPU-bound; and CARD remains a fallback/probe
path rather than faithful CARD reproduction~\cite{wei2026card}. D2 is not
causal collision harm; D3 uses prefix-neighborhood and reranker checks, not
trained-generator ranking; D5 is structural cost rather than measured latency;
D6 is optional churn evidence; and D7 requires generator traces absent from the
current rows.

New adapters emit the same normalized tables, declare their evidence role, and
record feature or checkpoint provenance. A named tokenizer improves coverage
only when item-level codes pass the validator and join to metadata or
interactions; a controlled mechanism probe improves diagnostic calibration but
does not count as named-method coverage; a generator trace activates D7.
Future versions should add causal collision-harm validation, broader faithful
adapter coverage, generator-output diagnostics, and unified search/recommendation
checks~\cite{li2024survey,penha2025jointsid}. Industrial SID deployment and
ranking studies further motivate treating these artifact properties as
engineering objects~\cite{singh2023bettersemanticids,li2026sidcoord,ju2026snapchatsid}.
